\documentclass[a4paper,12pt]{article}

\usepackage[utf8]{inputenc}
\usepackage[russian]{babel}
\usepackage{xstring}
\usepackage{amsmath}
\usepackage{amsfonts}
\usepackage{amssymb}
\usepackage{amsthm}
\usepackage{enumitem}
\usepackage{multicol}
\usepackage[svgnames]{xcolor}
\usepackage{sectsty}
\usepackage{soul}
\usepackage{soulutf8}
\usepackage{pb-diagram}
\usepackage{graphicx}
\usepackage{tikz}
\usepackage{rotating}

\author{
Типографические приемы
Выполнил: Фамилия И.\,О. \\
Группа: группа
}
\title{Лабораторная работа по \LaTeX}
\date{\today}

\sodef\soparam{}{.5ex plus .2ex}{1em plus .5 em}{1em plus .5 em}
\def\so#1{\edef\xso{{#1}}\expandafter\soparam\xso}

\voffset=-25.4mm
\topmargin=25mm
\headheight=0mm
\headsep=0mm
\textheight=247mm

\hoffset=-25.4mm
\oddsidemargin=20mm
\evensidemargin=20mm
\textwidth=170mm

\allsectionsfont{\raggedright\rmfamily}
\makeatletter
\def\@seccntformat#1{\csname the#1\endcsname.~}
\makeatother

\makeatletter
\def\eatrelax relax#1#2EndOfControl{#1#2} \relax
\def\createverb#1{%
    \begingroup
    \xdef\bbb{\detokenize{#1}}
    \let\do\@makeother
    \dospecials
    \catcode`\@=12
    \catcode`\{=12
    \catcode`\}=12
    \catcode`\\=12
    \catcode`\ =10
    \tokenize{\verbtext}{\bbb}%
    \ttfamily\expandafter\eatrelax\verbtext EndOfControl
    \endgroup
}%
\def\hmrg{.7em}
\def\vmrg{.7ex}
\def\frm#1#2#3{\fcolorbox{black}{#3}{\hspace*{\hmrg}\begin{minipage}{#1}\vspace*{\vmrg}#2\vspace*{\vmrg}\end{minipage}\hspace*{\hmrg}}}
\long\def\tsk#1{\fcolorbox{black}{MistyRose}{\hspace*{\hmrg}\begin{minipage}{.8\textwidth}#1\end{minipage}\hspace*{\hmrg}}}
\def\illustr#1#2#3{\par\vspace*{\vmrg}\noindent\frm{.55\textwidth}{\createverb{#1}}{#2}\quad \frm{.35\textwidth}{#1}{#3}\par\vspace*{\vmrg}}
\def\noillustr#1#2#3{\par\vspace*{\vmrg}\noindent\frm{.55\textwidth}{\createverb{#1}}{#2}\par\vspace*{\vmrg}}
\def\illustrate#1{\illustr{#1}{Moccasin}{Azure}}
\def\noillustrate#1{\noillustr{#1}{Moccasin}{Azure}}
\def\nillustrate#1{\illustr{#1}{yellow}{Pink}}
\makeatother

\newtheorem{thr}{Теорема}
\newtheorem{crlx}[thr]{Следствие}
\newtheorem{lmm}[thr]{Лемма}
\newtheorem{clm}[thr]{Утверждение}
\newtheorem{dfn}{Определение}
\newtheorem{ax}{Аксиома}

\newtheoremstyle{claim}{1mm}{.2mm}{\slshape}{}{\scshape}{:}{.5em plus 3em}{}
\theoremstyle{claim}
\newtheorem{crl}[thr]{Следствие}

\newtheoremstyle{theor}{1mm}{.2mm}{\upshape\sffamily}{}{\scshape}{:}{.5em}{}
\theoremstyle{theor}
\newtheorem{thrx}[thr]{Теорема}

\begin{document}

\maketitle

\textbf{12. Задания}
Написать в \LaTeX\ следующие фрагменты:

\vspace{1em}
\noindent 1)\hspace{0.7em}
\tsk{%
\begin{tabular}{p{0.48\linewidth}|p{0.48\linewidth}}
\textbf{Теорема 1.} \textit{Если $a$ зависит от $b$, то $b$ зависит от $a$.}

\textit{Доказательство.} Основано на использовании теоремы Лахлана об определимости.
Обычно, сначала доказательство проводится, &
когда рассматривается зависимость над $\omega$-насыщенной элементарной подсистемой, после чего результат обобщается на произвольный случай.
\end{tabular}
}

\vspace{1em}
\noindent 3)\hspace{0.7em}
\tsk{%
\textbf{Определение 1} (Ортогональная матрица). Матрица $A$ называется \so{ортогональной}, если $A^{-1}=A^{T}$.
}

\vspace{1em}
\noindent 4)\hspace{0.7em}
\tsk{%
Следующий граф нельзя обойти, пройдя по каждому из ребер
в точности один раз:

\vspace{0.8em}
\begin{center}
\begin{tikzpicture}[line cap=round, line join=round]
  \node (A) at (0,0) {$A$};
  \node (B) at (2.4,0) {$B$};
  \node (C) at (4.8,0) {$C$};
  \node (D) at (2.4,-1.8) {$D$};
  \draw (A) -- (B);
  \draw (B) -- (C);
  \draw (A) -- (D);
  \draw (B) -- (D);
  \draw (C) -- (D);
\end{tikzpicture}
\end{center}
}

\vspace{1em}
\noindent 5)\hspace{0.7em}
\tsk{%
Определение 2 (Амальгамируемость). Категория обладает
\so{свойством амальгамируемости}, если любая диаграмма морфизмов вида

\vspace{0.8em}
\begin{center}
$$\begin{diagram}
\node{A}\arrow{e}\arrow{s}\node{B}\\
\node{C}
\end{diagram}$$
\end{center}

может быть достроена до коммутативной:

\vspace{0.6em}
\begin{center}
$$\begin{diagram}
\node{A}\arrow{e}\arrow{s}\node{B}\arrow{s,..}\\
\node{C}\arrow{e,..}\node{D}
\end{diagram}$$
\end{center}
}

\vspace{1em}
\noindent 6)\hspace{0.7em}
\tsk{%
Гомоморфный образ $(\mathfrak{H})$ группы $(\mathfrak{G})$

По закону коммунизма

Изоморфен фактор-группе $(\mathfrak{G}/\mathfrak{K})$

По ядру гомоморфизма $(\mathfrak{K})$.

\vspace{0.8em}
\begin{center}
$$\begin{diagram}
\node{\mathfrak{G}}\arrow{e,t}{\supseteq}\arrow{se,t}{h}\node{\mathfrak{K}}\\
\node{\mathfrak{G}/\mathfrak{K}}\arrow{e,<>}\node{\mathfrak{H}}
\end{diagram}$$
\end{center}
}

\vspace{1em}
\noindent 7)\hspace{0.7em}
\tsk{%
Таблица смешения цветов:

\vspace{0.8em}
\begin{center}
\setlength{\tabcolsep}{10pt}
\renewcommand{\arraystretch}{1.55}
\begin{tabular}{|c|c|c|c|}
\hline
 & \rotatebox{90}{\textcolor{red}{Красный}} & \rotatebox{90}{\textcolor{green!70!black}{Зеленый}} & \rotatebox{90}{\textcolor{blue}{Синий}} \\
\hline
\textcolor{red}{Красный} & \textcolor{red}{\rule{1.6ex}{1.6ex}} & \textcolor{yellow}{\rule{1.6ex}{1.6ex}} & \textcolor{magenta}{\rule{1.6ex}{1.6ex}} \\
\hline
\textcolor{green!70!black}{Зеленый} & \textcolor{yellow}{\rule{1.6ex}{1.6ex}} & \textcolor{green!70!black}{\rule{1.6ex}{1.6ex}} & \textcolor{cyan}{\rule{1.6ex}{1.6ex}} \\
\hline
\textcolor{blue}{Синий} & \textcolor{magenta}{\rule{1.6ex}{1.6ex}} & \textcolor{cyan}{\rule{1.6ex}{1.6ex}} & \textcolor{blue}{\rule{1.6ex}{1.6ex}} \\
\hline
\end{tabular}
\end{center}
}

\end{document}
